% answers_short.tex —— 2019 年·全国Ⅱ卷·理科 · 简短答案速查
% 由 tools/make_short_answers.py 自 answers.tex 的 \gksolhead/\gkanswer 提取，
% 并与 2019/2019_数学真题_全国2理_答案要点.md 逐题比对后生成，勿手改（改源 MD 后重跑）。

\begin{gknotice}
\textbf{使用说明}：本节只给最终答案便于快速核对，完整推导见解析版。答案为编写者独立推导，\textbf{非官方标准答案}。
\end{gknotice}

\gkpart{选择题}
\begin{center}\begin{tabular}{ccccccccccccc}
\toprule
题号 & 1 & 2 & 3 & 4 & 5 & 6 & 7 & 8 & 9 & 10 & 11 & 12 \\
\midrule
答案 & A & C & C & D & A & C & B & D & A & B & A & B \\
\bottomrule
\end{tabular}\end{center}

\gkpart{填空题}
\begin{itemize}[leftmargin=2.6em]
  \item[\textbf{13.}] $0.98$
  \item[\textbf{14.}] $-3$
  \item[\textbf{15.}] $6\sqrt{3}$
  \item[\textbf{16.}] $26$；$\sqrt{2}-1$
\end{itemize}

\gkpart{解答题}
\begin{itemize}[leftmargin=2.6em]
  \item[\textbf{17.}] (1) 证明见解析；(2) $\dfrac{\sqrt{3}}{2}$
  \item[\textbf{18.}] (1) $0.5$；(2) $0.1$
  \item[\textbf{19.}] (1) 证明见解析；(2) $a_n = \left(\dfrac{1}{2}\right)^n + n - \dfrac{1}{2}$，$b_n = \left(\dfrac{1}{2}\right)^n - n + \dfrac{1}{2}$
  \item[\textbf{20.}] (1) 单调递增区间为 $(0, 1)$ 和 $(1, +\infty)$；有且仅有两个零点；(2) 证明见解析
  \item[\textbf{21.}] (1) $\dfrac{x^2}{4} + \dfrac{y^2}{2} = 1$ ($x \neq \pm 2$)；曲线 $C$ 为椭圆（除去左右顶点）；(2) ① 证明见解析；② $\dfrac{16}{9}$
\end{itemize}

\gkpart{选做题}
\begin{itemize}[leftmargin=2.6em]
  \item[\textbf{22.}] (1) $\rho_0 = 2\sqrt{3}$；$l$ 的极坐标方程为 $\rho\cos\left(\theta - \dfrac{\pi}{3}\right) = 2$（或 $\rho\sin\left(\theta + \dfrac{\pi}{6}\right) = 2$）；(2) $\rho = 4\cos\theta$ $\left(\dfrac{\pi}{4} \leqslant \theta \leqslant \dfrac{\pi}{2}\right)$
  \item[\textbf{23.}] (1) $\{x \mid x < 1\}$；(2) 按原卷字面条件无解（$\varnothing$）；若勘误为开区间 $x \in (-\infty, 1)$ 则 $a$ 的取值范围为 $[1, +\infty)$
\end{itemize}
